% =============================================================================
% AI Tageszeitung — Rheinisch-Bergischer Kreis — TITELSEITE (Seite 1)
% Kreisweite Übersicht: Top-Nachrichten aller 8 Kommunen
% Kompilierung: lualatex titelseite.tex
% =============================================================================
\documentclass[11pt, a4paper]{scrartcl}

% --- Pakete ---
\usepackage{fontspec}
\usepackage[ngerman]{babel}
\usepackage[top=12mm, bottom=15mm, left=10mm, right=10mm]{geometry}
\usepackage{multicol}
\usepackage{graphicx}
\usepackage{xcolor}
\usepackage{hyperref}
\usepackage{tikz}
\usepackage{tcolorbox}
\usepackage{enumitem}
\usepackage{datetime2}

% --- Farben (Zeitungs-Design) ---
\definecolor{rbkblau}{HTML}{1B3A5C}
\definecolor{rbkrot}{HTML}{C0392B}
\definecolor{rbkgrau}{HTML}{F5F5F5}
\definecolor{akzent}{HTML}{2980B9}
\definecolor{liniengrau}{HTML}{BFBFBF}

% --- Schriften ---
\setmainfont{TeX Gyre Termes}        % Serif für Fließtext
\setsansfont{TeX Gyre Heros}         % Sans für Headlines
\newfontfamily\headlinefont{TeX Gyre Heros Bold}

% --- Hyperlinks ---
\hypersetup{
  colorlinks=true,
  linkcolor=rbkblau,
  urlcolor=akzent,
  pdftitle={AI Tageszeitung -- Rheinisch-Bergischer Kreis},
}

% --- Makros ---
\newcommand{\zeitstempel}{\DTMnow}
\newcommand{\topnachricht}[4]{%
  % #1=Kommune, #2=Headline, #3=Teaser, #4=Link
  \begin{tcolorbox}[
    colback=white,
    colframe=liniengrau,
    left=3mm, right=3mm, top=2mm, bottom=2mm,
    boxrule=0.3pt,
    arc=1mm
  ]
  {\small\sffamily\color{rbkrot}\textbf{#1}}\\[1pt]
  {\large\headlinefont #2}\\[3pt]
  {\small #3}\\[2pt]
  {\footnotesize\color{akzent}\href{#4}{Weiterlesen \textrightarrow}}
  \end{tcolorbox}
  \vspace{1mm}
}

\newcommand{\tickeritem}[2]{%
  % #1=Zeit, #2=Kurzmeldung
  \item[\textbf{\small\color{rbkblau}#1}] {\small #2}
}

% --- Dokument ---
\begin{document}
\pagestyle{empty}

% === KOPFZEILE / MASTHEAD ===
\begin{center}
\begin{tikzpicture}
  \fill[rbkblau] (0,0) rectangle (\textwidth, 1.8cm);
  \node[anchor=west, white] at (0.3, 0.9) {%
    {\fontsize{28}{30}\selectfont\headlinefont AI TAGESZEITUNG}};
  \node[anchor=east, white] at (\textwidth-0.3, 1.2) {%
    {\small\sffamily Rheinisch-Bergischer Kreis}};
  \node[anchor=east, white] at (\textwidth-0.3, 0.5) {%
    {\footnotesize\sffamily Aktualisiert: \zeitstempel}};
\end{tikzpicture}
\end{center}

\vspace{2mm}

% === LAUFBAND / TICKER ===
\begin{tcolorbox}[
  colback=rbkgrau,
  colframe=rbkblau,
  boxrule=0.5pt,
  left=3mm, right=3mm, top=1mm, bottom=1mm
]
{\sffamily\small\textbf{EILMELDUNGEN:}}\quad
%%TICKER_PLACEHOLDER%%
\end{tcolorbox}

\vspace{3mm}

% === HAUPTBEREICH: 3-Spalten-Layout ===
\begin{multicols}{3}
\raggedcolumns

% --- Aufmacher (spaltenübergreifend wird via minipage gelöst) ---
\section*{}

%%TOPNACHRICHTEN_PLACEHOLDER%%

% Platzhalter-Beispiele:
\topnachricht{Bergisch Gladbach}
  {Neue S-Bahn-Haltestelle beschlossen}
  {Der Rat hat einstimmig den Bau einer neuen Haltestelle am Technologiepark genehmigt. Baubeginn ist für 2027 geplant.}
  {https://example.com/artikel/1}

\topnachricht{Overath}
  {Hochwasserschutz: Agger-Deich wird erhöht}
  {Nach den Überschwemmungen 2021 investiert die Stadt 4,2 Mio. Euro in präventive Maßnahmen.}
  {https://example.com/artikel/2}

\topnachricht{Kürten}
  {Glasfaser-Ausbau erreicht alle Ortsteile}
  {Bis Ende 2026 sollen auch Dürscheid und Biesfeld vollständig angeschlossen sein.}
  {https://example.com/artikel/3}

\end{multicols}

% === FUßZEILE ===
\vfill
\begin{center}
\rule{\textwidth}{0.4pt}\\[2mm]
{\footnotesize\sffamily\color{rbkblau}
  Bergisch Gladbach~|~Burscheid~|~Kürten~|~Leichlingen~|~Odenthal~|~Overath~|~Rösrath~|~Wermelskirchen\\[1mm]
  Automatisch aktualisiert aus lokalen RSS-Quellen~|~Keine Gewähr für Vollständigkeit\\[1mm]
  \href{https://github.com/dieternadolski/ai-tageszeitung}{GitHub}~|~Kontakt: dieter.nadolski@gmail.com
}
\end{center}

\end{document}
